\section{Background}
\label{sec:background}

\subsection{Threat model}
The attacker controls a share of the training data and plants a trigger with a target label. The defender receives the trained model and holds a small clean validation set drawn from the same distribution, here 2000 images, but never the training set and never a poisoned example. At deployment the defender sees 1 input at a time and must decide whether it carries a trigger before returning a prediction. This is the setting of STRIP \cite{gao2019strip}, SCALE-UP \cite{guo2023scale}, IBD-PSC \cite{hou2024ibdpsc}, TeCo \cite{liu2023teco} and PSBD, and it excludes detectors that need the poisoned training pool or that retrain the model.

\subsection{Prediction shift and PSBD}
PSBD \cite{li2025psbd} runs the model $k$ times on an input $x$ with dropout at rate $p$ active and compares each dropout prediction with the prediction without dropout. The paper defines the prediction shift indicator and the shift ratio of a set $\mathcal{D}$ as
\begin{equation}
  \ps(x) = \mathbb{I}\big[\mathcal{Y}(x;\theta) \neq \mathcal{Y}(x;\theta')\big], \qquad
  \sigma(\mathcal{D}) = \frac{1}{k|\mathcal{D}|}\sum_{x\in\mathcal{D}} \ps(x),
\end{equation}
where $\theta$ are the weights, $\theta'$ the weights with dropout active and $\mathcal{Y}$ the predicted class. The detection statistic, the prediction shift uncertainty, is the drop in the probability of the no-dropout class averaged over the $k$ passes,
\begin{equation}
  \psu(x) = \frac{1}{k}\sum_{i=1}^{k} \big[ P(\hat y \mid x;\theta) - P(\hat y \mid x;\theta'_i) \big],
\end{equation}
with $\hat y$ the class predicted without dropout. A poisoned input keeps its probability under dropout and has a low $\psu$. The dropout rate is chosen so that the shift ratio of the clean validation set reaches 0.8, the threshold is the 25th percentile of the validation set's $\psu$, and $k = 3$. The paper places dropout after each residual connection of a ResNet basic block and explains the phenomenon by neuron bias: dropout removes clean features, the network falls back on the classes its neurons are biased toward, and the backdoor's neurons are biased more firmly than the rest. We return to that explanation in \cref{sec:mechanism}.

\subsection{Where a perturbation can go in a Vision Transformer}
A ViT-B/16 \cite{dosovitskiy2021an} splits a 224 by 224 image into 196 patches of 16 by 16 pixels, embeds each as a 768-dimensional token, prepends a class token and runs 12 identical blocks. Each block applies a LayerNorm, multi-head self-attention over the 197 tokens and a residual add, then a second LayerNorm, a 2-layer MLP applied to every token independently and a second residual add. The classifier reads the class token after a final LayerNorm. \Cref{fig:vit-block} shows 1 block with the sites where we inject a perturbation. Attention is the only operation that mixes information across tokens, so a perturbation before the attention LayerNorm acts on each token before any mixing, a perturbation before the MLP LayerNorm acts after the block's mixing has happened, and a perturbation on the residual stream acts on the sum of everything the network has written so far. Swin-S \cite{liu2021swin} keeps the block structure but restricts attention to shifted windows of 7 by 7 tokens and merges patches between its 4 stages, so tokens mix locally within a stage and globally only through merging.

\begin{figure}[t]
\centering
\begin{tikzpicture}[
  node distance=3mm,
  box/.style={draw, rounded corners=1pt, minimum width=22mm, minimum height=6mm, font=\scriptsize, align=center},
  norm/.style={box, fill=gray!12},
  op/.style={box, fill=blue!8},
  plus/.style={draw, circle, inner sep=1pt, font=\scriptsize},
  site/.style={draw, circle, fill=orange!80, inner sep=1.4pt, font=\bfseries\tiny},
  flow/.style={-{Latex[length=1.6mm]}, thick},
  lbl/.style={font=\scriptsize, text width=34mm, align=left}
]
\node[font=\scriptsize] (in) {residual stream $x$ \quad (197 tokens $\times$ 768)};
\node[norm, below=of in] (ln1) {LayerNorm};
\node[op, below=of ln1] (attn) {multi-head attention};
\node[plus, below=of attn] (add1) {$+$};
\node[norm, below=of add1] (ln2) {LayerNorm};
\node[op, below=of ln2] (mlp) {MLP (per token)};
\node[plus, below=of mlp] (add2) {$+$};
\node[font=\scriptsize, below=of add2] (out) {to the next block};
\draw[flow] (in) -- (ln1); \draw[flow] (ln1) -- (attn); \draw[flow] (attn) -- (add1);
\draw[flow] (add1) -- (ln2); \draw[flow] (ln2) -- (mlp); \draw[flow] (mlp) -- (add2); \draw[flow] (add2) -- (out);
\draw[flow] (in.south) ++(0,-1mm) -| ++(-16mm,0) |- (add1.west);
\draw[flow] (add1.south) ++(0,-1mm) -| ++(-16mm,0) |- (add2.west);
\node[site, right=2mm of ln1.north east, yshift=1mm] (A) {A};
\node[site, right=2mm of ln2.north east, yshift=1mm] (B) {B};
\node[site, right=2mm of attn.south east, yshift=-1mm] (C) {C};
\node[site, right=2mm of add1.east, xshift=4mm] (D) {D};
\node[lbl, right=10mm of A] {A \textbf{before the attention norm}: token mask (recommended), channel mask, Gaussian noise};
\node[lbl, right=10mm of C] {C \textbf{before the residual add}: dropout on the branch output};
\node[lbl, right=10mm of D] {D \textbf{after the residual add}: dropout on the stream (the published placement)};
\node[lbl, right=10mm of B] {B \textbf{before the MLP norm}: token mask, Gaussian noise};
\end{tikzpicture}

\caption{One ViT block with the positions we perturb. Token masking at the input of the attention LayerNorm (A) is the recommended placement. Dropout on the residual stream after the add (D) is the placement the PSBD paper uses on ResNets. The MLP input (B), the branch output before the add (C), the attention heads and the MLP hidden units are the other sites in the basis.}
\label{fig:vit-block}
\end{figure}
